\documentclass[11pt,a4paper]{article}
\usepackage[T1]{fontenc}
\usepackage[utf8]{inputenc}
\usepackage[main=italian,provide=*]{babel}
\usepackage{amsmath,amssymb}
\usepackage{geometry}
\geometry{margin=2.3cm,top=2.5cm,bottom=2.7cm}
\usepackage{booktabs}
\usepackage[table]{xcolor}
\usepackage{tcolorbox}
\tcbuselibrary{skins,breakable}
\usepackage{titlesec}
\usepackage{enumitem}
\usepackage{seqsplit}
\usepackage{needspace}
\usepackage{fancyhdr}
\usepackage{hyperref}
\hypersetup{colorlinks=true,linkcolor=Sepia,citecolor=Sepia,urlcolor=Sepia}

\definecolor{inkblue}{HTML}{1B3A5C}
\definecolor{lightblue}{HTML}{EAF1F8}
\definecolor{accent}{HTML}{B0562A}
\definecolor{softgray}{HTML}{6B6B6B}
\definecolor{okgreen}{HTML}{2E6B3E}
\definecolor{okgreenbg}{HTML}{EAF5EC}
\definecolor{warnamber}{HTML}{9C6B12}
\definecolor{warnbg}{HTML}{FBF1E0}
\definecolor{advisorpurple}{HTML}{5B3A7A}
\definecolor{advisorbg}{HTML}{F1ECF6}
\definecolor{notebar}{HTML}{9C6B12}

\titleformat{\section}
  {\normalfont\Large\bfseries\color{inkblue}}
  {\thesection}{0.6em}{}[{\color{inkblue!40}\titlerule}]
\titlespacing{\section}{0pt}{0.8em}{0.4em}
\titleformat{\subsection}
  {\normalfont\large\bfseries\color{accent}}
  {}{0em}{}
\titlespacing{\subsection}{0pt}{0.5em}{0.2em}

\pagestyle{fancy}
\fancyhf{}
\renewcommand{\headrulewidth}{0.4pt}
\fancyhead[R]{\small\color{softgray}\thepage}
\fancyfoot[C]{\small\color{softgray}Tesi triennale, Samuele Schiavi --- Relatore: Prof.\ Alessandro Chiesa}

\newtcolorbox{keyeq}[1][]{
  colback=lightblue, colframe=inkblue!70, boxrule=0.6pt,
  arc=2pt, left=8pt, right=8pt, top=4pt, bottom=4pt, #1
}
\newtcolorbox{panelbox}[1]{
  colback=white, colframe=softgray!50, boxrule=0.5pt,
  arc=2pt, left=7pt, right=7pt, top=3pt, bottom=3pt,
  fonttitle=\bfseries\color{inkblue}, title={#1}
}
\newtcolorbox{okbox}[1]{
  colback=okgreenbg, colframe=okgreen!60, boxrule=0.6pt,
  arc=2pt, left=7pt, right=7pt, top=3pt, bottom=3pt,
  fonttitle=\bfseries\color{okgreen}, title={#1}
}
\newtcolorbox{warnbox}[1]{
  colback=warnbg, colframe=warnamber!70, boxrule=0.6pt,
  arc=2pt, left=7pt, right=7pt, top=3pt, bottom=3pt,
  fonttitle=\bfseries\color{warnamber}, title={#1}
}
\newtcolorbox{advisorbox}[1]{
  colback=advisorbg, colframe=advisorpurple!60, boxrule=0.6pt,
  arc=2pt, left=7pt, right=7pt, top=4pt, bottom=4pt,
  fonttitle=\bfseries\color{advisorpurple}, title={#1}
}
\newtcolorbox{editnote}{
  colback=white, colframe=white, boxrule=0pt,
  borderline west={2.2pt}{0pt}{notebar},
  arc=0pt, left=10pt, right=4pt, top=2pt, bottom=2pt,
  fontupper=\itshape\small\color{softgray!90!black}
}
\fancyhead[L]{\small\color{softgray}Circuito compatto per il passo di Trotter --- trimero catena aperta}

\title{\vspace{-1.5em}\color{inkblue}Circuito compatto per il Trotter del trimero a catena aperta\\[0.1em]
\Large\normalfont\color{softgray}Verifica empirica, mirror diretto dello studio già fatto per l'anello}
\author{Note di lavoro --- Tesi triennale, Samuele Schiavi\\ Relatore: Prof.\ Alessandro Chiesa}
\date{}

\begin{document}
\sloppy
\maketitle
\thispagestyle{fancy}
\vspace{-2em}

Seguito naturale di \texttt{trimero\_anello\_circuito\_compatto.pdf}, applicato
ora alla catena aperta. Stessa domanda, stesso metodo, punto di lavoro
diverso ($S_0$ invece di $R_0$). Riportato per completezza, non perché ci
si aspettasse un esito qualitativamente diverso.

\section{Il problema, invariato rispetto all'anello}

\begin{warnbox}{Stessa differenza dal dimero, qui confermata}
Per il trimero (3 qubit, $SU(8)$), anello o catena che sia, non esiste un
teorema di compressione universale chiuso come quello a 2 qubit del
dimero (Cartan/Weyl, Vatan--Williams 2004). Anche qui la verifica è
\textbf{empirica} (transpiler), non un teorema.
\end{warnbox}

Il circuito Trotter della catena (\texttt{trotter\_trimero\_catena.py}) usa,
per un singolo passo esterno al punto $S_0$ ($J{=}1,b{=}3.0073414,D{=}0.3$),
\textbf{10 gate a due qubit nativi} ($2\times R_{XX}+2\times R_{YY}+6\times
R_{ZZ}$: due bond per lo scambio, due bond per il DM), più 8 Hadamard e 3
$R_z$ per il campo -- meno della metà dei 15 dell'anello, proporzionale al
numero di bond (2 contro 3).

\section{Metodo, identico all'anello}

\begin{enumerate}[leftmargin=1.4em,itemsep=0.25em]
\item costruzione del circuito Trotter esplicito con
\texttt{trotter\_trimero\_catena.py};
\item estrazione dell'unitaria $8\times8$ esatta, \texttt{Operator(qc).data}
(già validata a precisione macchina contro la matrice indipendente, self-test
5 del modulo);
\item la matrice passata come blocco opaco, \texttt{qc.unitary(U,[0,1,2])};
\item sintesi: \texttt{transpile(qc, basis\_gates=['cx','rz','sx','x'],
optimization\_level=3)};
\item verifica di fedeltà (\texttt{process\_fidelity}) fra compresso ed
esplicito.
\end{enumerate}

\section{Risultati}

\subsection*{Un solo passo, a $S_0$}

\begin{center}
\renewcommand{\arraystretch}{1.3}
\begin{tabular}{lcc}
\toprule
& gate a due qubit & fedeltà vs esplicito \\
\midrule
circuito fisico (esplicito) & 10 ($R_{XX}{+}R_{YY}{+}R_{ZZ}$ nativi) & --- \\
compresso (\texttt{optimization\_level=3}) & \textbf{19 CNOT} & $0.9999999999999996$ \\
\bottomrule
\end{tabular}
\end{center}

Fedeltà a precisione macchina (scarto $\sim4\times10^{-16}$, non zero
esatto per via degli arrotondamenti del transpiler stesso, non un errore
di approssimazione fisica). Il compresso usa \emph{più} gate a due qubit
del fisico esplicito qui (19 contro 10): non contraddittorio, il conteggio
di riferimento rilevante è quello del prodotto a molti passi (sotto), non
del singolo passo isolato -- per un solo passo il circuito esplicito è già
economico, poco spazio per il transpiler di comprimere ulteriormente.

\subsection*{L'intero prodotto a $N$ passi --- il punto concettualmente importante}

\begin{center}
\renewcommand{\arraystretch}{1.3}
\begin{tabular}{cccc}
\toprule
\rowcolor{inkblue!10}
$N$ & gate 2q nativi (esplicito) & CNOT (compresso) & fedeltà \\
\midrule
1   & 10   & 19 & $1.0$ \\
2   & 20   & 19 & $1.0$ \\
5   & 50   & 19 & $1.0$ \\
20  & 200  & 19 & $1.0$ \\
100 & 1000 & 19 & $1.0$ \\
\bottomrule
\end{tabular}
\end{center}

Il conteggio compresso si stabilizza a $19$ CNOT già da $N=1$ in poi,
verificato fino a $N=100$: identico fenomeno di anello e dimero. A
$N=8000$ (soglia di convergenza usata nello studio principale, Sezione~6 di
\texttt{\seqsplit{quantum\_simulation\_trimero\_catena\_applicazione.tex}}), il
circuito esplicito avrebbe $80000$ gate a due qubit, contro i $\sim19$ del
compresso.

\section{Perché non si usa per il rumore (stessa conclusione di anello e dimero)}

\begin{warnbox}{Il compromesso $N$-vs-rumore sparirebbe}
Identica considerazione già fatta per anello e dimero: il costo in gate del
compresso resta fisso a $\sim19$ qualunque sia $N$, quindi userlo in Parte 2
farebbe sparire dal grafico il compromesso fra errore di Trotter
($\propto1/N$) ed errore fisico di gate sotto rumore (scala col numero di
gate). Con gli $N$ dell'ordine delle migliaia trovati per $S_0$, l'effetto
sarebbe ancora più ingannevole: si misurerebbe il rumore su un'unitaria a
costo fisico del tutto slegato dal numero di passi realmente usati per la
fisica.
\end{warnbox}

\section{Riepilogo}

\begin{enumerate}[leftmargin=1.4em,itemsep=0.3em]
\item Nessun teorema di compressione universale per 3 qubit: verifica
empirica, non un minimo dimostrato -- stessa situazione dell'anello.
\item Un singolo passo (10 gate nativi a $S_0$) si comprime a 19 CNOT,
fedeltà $\approx1.0$.
\item Il prodotto a $N$ passi si comprime a $\sim19$ CNOT
\textbf{indipendentemente da $N$} (verificato fino a $N=100$) -- stesso
fenomeno di anello e dimero, qui con un rapporto di compressione ancora
più marcato per la scala di $N$ richiesta ($\sim8000$ contro $\sim1000$
dell'anello).
\item Conclusione invariata: non usare il circuito compresso per lo studio
Trotter-vs-rumore della Parte 2.
\end{enumerate}

\needspace{15\baselineskip}
\section*{Appendice: comandi Qiskit usati per la verifica}
\addcontentsline{toc}{section}{Appendice}

\begin{panelbox}{Codice}
\small
\begin{verbatim}
from qiskit import QuantumCircuit, transpile
from qiskit.quantum_info import Operator, process_fidelity
from trotter_trimero_catena import trotter_circuit

qc_step = trotter_circuit(J, b, D, tau, 1)          # un passo esplicito
U = Operator(qc_step).data                          # unitaria 8x8 esatta

qc_compact = QuantumCircuit(3)
qc_compact.unitary(U, [0, 1, 2])                    # blocco opaco
qc_t = transpile(qc_compact, basis_gates=['cx','rz','sx','x'],
                  optimization_level=3)

print(qc_t.count_ops())
print(process_fidelity(Operator(qc_t), Operator(qc_step)))

# ripetuto su trotter_circuit(J, b, D, T, N) per N = 1, 2, 5, 20, 100
\end{verbatim}
\end{panelbox}

\end{document}
